\section{Empirical Demonstrations}
\label{sec:experiments}

The preceding sections develop the Bayesian/Kolmogorov framework
theoretically. This section demonstrates it empirically. Five experiments,
each isolating one aspect of the objective function $L(H) + L(D \mid H)$,
run as code with reproducible numbers. All results are available in the
\texttt{auto-research} repository; the command to reproduce each experiment
is noted.

These are not five different demonstrations. They are five windows into the
same object. \texttt{compress} measures $L(H)$ directly.
\texttt{mdl} shows the tradeoff between $L(H)$ and $L(D \mid H)$ in model
selection. \texttt{replicate} shows what happens when $L(H)$ is
underreported. \texttt{search} shows that verification is just evaluating
$L(D \mid H)$ cheaply, and that this cheapness is what makes large-scale
search tractable. \texttt{update} shows the Kolmogorov prior doing work
that data alone cannot do. The sixth experiment, \texttt{autoloop}, brings
all five together in a running AutoResearch loop and identifies the precise
gap between current systems and the full MDL objective.

\subsection{Structure Is Compressibility (\texttt{ar compress})}

Kolmogorov complexity is uncomputable, but lossless compression provides
a computable upper bound. A string that compresses well has low $\K$; one
that does not compress has high $\K$. The compression ratio approximates
$\K(x)/|x|$---the complexity relative to the raw length.

Running \texttt{ar compress} on files with known structural properties
yields:

\begin{center}
\begin{tabular}{lrrl}
\toprule
File & Raw bytes & Ratio & Interpretation \\
\midrule
\texttt{repeated\_pattern} & 1009 & 0.036 & highly structured \\
\texttt{genome\_like}      &  766 & 0.123 & highly structured \\
\texttt{source\_code}      &  868 & 0.414 & moderate structure \\
\texttt{english\_prose}    & 1501 & 0.502 & moderate structure \\
\texttt{pi\_digits}        & 1009 & 0.520 & moderate structure \\
\texttt{random\_noise}     &  369 & 0.564 & low structure \\
\texttt{natural\_law}      &  591 & 0.604 & low structure \\
\bottomrule
\end{tabular}
\end{center}

Several results deserve comment.

The genome compresses to 12.3\% of its raw size. DNA looks complex---a
four-character alphabet, long sequences, no obvious pattern to the eye. But
it is built from repeated motifs, and evolution is a compressive process.
Biological structure is compressible structure.

The digits of $\pi$ compress to only 52.0\%, which surprises most people.
$\pi$ has very low true Kolmogorov complexity: there exists a short program
(any standard $\pi$-computation algorithm) that generates it exactly. But
gzip cannot find that program---it finds only local repetitions. This
illustrates the gap between the compression proxy and the true $\K(x)$:
gzip is an upper bound, not the true complexity.

Natural laws compress to 60.4\%, below random noise. The equations use
diverse symbols with little local repetition; gzip finds structure by
matching substrings, and the structure in natural laws is semantic, not
syntactic. The true $\K$ of Newton's law---the length of the shortest
program that predicts all observations it covers---is extremely low. The
compression proxy cannot see it.

\subsection{The MDL Tradeoff Is Real (\texttt{ar mdl})}

MDL model selection fits polynomials of degree 1 through 8 to data generated
by a polynomial of known degree, computing $L(H) + L(D \mid H)$ for each.
The degree with the lowest total score is the MDL winner.

On data generated by a true degree-2 parabola with 40 points and noise
$\sigma = 0.3$:

\begin{center}
\begin{tabular}{rrrr}
\toprule
Degree & $L(H)$ & $L(D \mid H)$ & Total \\
\midrule
1 & 64.0 & 101.2 & 165.2 \\
2 & 96.0 & $-0.2$ & \textbf{95.8} \\
3 & 128.0 & $-0.4$ & 127.6 \\
4 & 160.0 & $-0.4$ & 159.6 \\
\bottomrule
\end{tabular}
\end{center}

MDL recovers the true degree. The jump in $L(D \mid H)$ from degree 1 to
degree 2 is 101.4 bits; degrees 3 and above improve fit negligibly while
$L(H)$ climbs linearly.

The decisive test is \texttt{overfit\_trap}: the same true degree-2
generator, but only 12 noisy points ($\sigma = 1.2$):

\begin{center}
\begin{tabular}{rrrr}
\toprule
Degree & $L(H)$ & $L(D \mid H)$ & Total \\
\midrule
1 & 64.0 & 31.4 & \textbf{95.4} \\
2 & 96.0 & 25.4 & 121.4 \\
3 & 128.0 & 24.9 & 152.9 \\
\bottomrule
\end{tabular}
\end{center}

MDL chooses degree 1, not the true degree 2. This is not a failure. With
12 noisy points, the improvement in fit from degree 2 does not justify the
32-bit complexity penalty. The framework is being epistemically honest: we
do not have enough data to establish the true degree. The simpler model is
the correct answer given the available evidence. Science should say the
same thing.

Note that $L(D \mid H)$ can be negative. When the model fits well, the
Gaussian log-likelihood is positive, making $-\log P(D \mid H)$ negative.
Negative $L(D \mid H)$ means the data is being compressed relative to the
coding baseline---the hypothesis is genuinely explaining structure.

\subsection{The Replication Crisis Is One Missing Term (\texttt{ar replicate})}

The \texttt{replicate} experiment simulates p-hacking on null data---pure
noise, no true effect---and computes the honest MDL correction alongside
the naive p-value.

\begin{center}
\begin{tabular}{rrrrl}
\toprule
Tests $N$ & Evidence (bits) & Search cost $\log_2 N$ & Corrected & Verdict \\
\midrule
10  & 5.13 & 3.32 & 1.80  & not significant \\
50  & 7.28 & 5.64 & 1.63  & not significant \\
100 & 6.28 & 6.64 & $-0.37$ & not significant \\
\bottomrule
\end{tabular}
\end{center}

Every naive result clears the 4.32-bit significance threshold
($p < 0.05$). Every corrected result is noise. At $N = 100$ the corrected
value is negative: the data provides less than zero evidence for the
hypothesis once the search cost is accounted for.

A single test on data with a real effect of size 0.8:

\begin{center}
\begin{tabular}{rrrrl}
\toprule
Tests $N$ & Evidence (bits) & Search cost $\log_2 N$ & Corrected & Verdict \\
\midrule
1   & 16.45 & 0.00 & 16.45 & significant \\
100 & 27.90 & 6.64 & 21.25 & significant \\
\bottomrule
\end{tabular}
\end{center}

A real effect survives correction with massive margin. The honest MDL
calculation does not suppress real discoveries; it suppresses artifacts.

The search cost $\log_2 N$ is exactly the Bonferroni correction, expressed
in bits. The statistics community had the right tool. The replication crisis
is the empirical consequence of not using it.

\subsection{Verification Is Constant; Search Scales (\texttt{ar search})}

The \texttt{search} experiment runs a pool of candidate implementations
against a locked test suite, measuring verification time per candidate.

\begin{center}
\begin{tabular}{lrr}
\toprule
Candidate & Tests passed & Verify time (ms) \\
\midrule
\texttt{broken\_empty\_return}  & 0/8 & 0.13 \\
\texttt{broken\_mutates\_input} & 0/8 & 0.62 \\
\texttt{broken\_off\_by\_one}   & 0/8 & 0.09 \\
\texttt{correct\_bubble}        & 8/8 & 0.07 \\
\texttt{correct\_recursive}     & 8/8 & 0.08 \\
\bottomrule
\end{tabular}
\end{center}

Verification time is flat across all candidates. Complexity of the
hypothesis does not affect verification cost. Correctness does not affect
verification cost. Search cost accumulates: 5 candidates $\times$ 0.16\,ms
$= 0.97$\,ms before the first passing candidate, a ratio of 6.1$\times$.
With 500 candidates the ratio would be ${\sim}600\times$; with 50,000,
${\sim}60{,}000\times$. The verification cost never moves.

The untrusted source consequence is concrete. The candidate
\texttt{broken\_returns\_none} prints to stdout, runs without error,
produces output, and returns \texttt{None}. The test suite catches this in
0.07\,ms, 0 of 12 tests passed. The verifier needed no knowledge of who
wrote the candidate or what they intended. The object spoke for itself.

\subsection{The Right Prior Converges Faster (\texttt{ar update})}

The \texttt{occams\_razor} case is the sharpest result. Four
hypotheses---$3x$, $3x + 0{\cdot}x^2$, $3x + 0{\cdot}x^3$,
$3x + 0{\cdot}x^2 + 0{\cdot}x^3$---predict identically for every input.

\begin{center}
\begin{tabular}{lrr}
\toprule
Step & Uniform on $3x$ & Kolmogorov on $3x$ \\
\midrule
Prior & 25.0\% & 99.6\% \\
1     & 25.0\% & 99.6\% \\
5     & 25.0\% & 99.6\% \\
\bottomrule
\end{tabular}
\end{center}

The uniform prior stays at 25\% forever. No finite amount of data can break
the tie when all likelihoods are identical. The Kolmogorov prior starts at
99.6\% on the simplest hypothesis before any data arrives. Among hypotheses
that predict equally well, the prior decides---and the prior should be the
Kolmogorov prior.

The \texttt{weak\_signal} case shows the indirect advantage. True hypothesis
$2x$, noise $= 1.5$:

\begin{center}
\begin{tabular}{lrr}
\toprule
Step & Uniform on $2x$ & Kolmogorov on $2x$ \\
\midrule
Prior & 25.0\% & 17.4\% \\
2     & 33.9\% & 33.7\% \\
3     & 68.5\% & 78.0\% \\
5     & 96.7\% & 98.9\% \\
7     & 99.9\% & 100.0\% \\
\bottomrule
\end{tabular}
\end{center}

The Kolmogorov prior starts $2x$ at 17.4\%---lower than uniform's 25\%,
because $x$ is a shorter expression and absorbs more prior mass. Yet by
step 3 the Kolmogorov posterior is at 78\% versus 68.5\% for uniform.
The mechanism: the Kolmogorov prior suppressed the wrong alternatives from
the start. When the likelihood accumulated on $2x$, there were fewer
competitors to absorb probability mass. Faster convergence means fewer
experiments needed to reach the same confidence---a compounding advantage
in any iterative research loop.

\subsection{The Karpathy Gap (\texttt{ar autoloop})}

The \texttt{autoloop} experiment models the Karpathy AutoResearch loop with
explicit MDL accounting, run in two modes: full MDL ($L(H) + L(D \mid H)$)
and $L(D \mid H)$ only---the current Karpathy objective.

On the \texttt{sort} problem, starting from a baseline that already passes
all 8 tests:

\begin{center}
\begin{tabular}{llrrrrl}
\toprule
Iter & Mutation & Tests & $L(H)$ & $L(D|H)$ & MDL & Decision \\
\midrule
0 & baseline               & 8/8 & 279 &  0.0 & 279.0 & --- \\
1 & \texttt{01\_fix\_break}        & 8/8 & 257 &  0.0 & 257.0 & keep \\
2 & \texttt{02\_add\_type\_check}  & 8/8 & 315 &  0.0 & 315.0 & discard \\
3 & \texttt{03\_simplify}          & 8/8 & 170 &  0.0 & 170.0 & keep \\
4 & \texttt{04\_introduce\_bug}    & 3/8 & 197 & 15.0 & 212.0 & discard \\
5 & \texttt{05\_use\_builtin}      & 8/8 &  97 &  0.0 &  97.0 & keep \\
6 & \texttt{06\_overengineered}    & 8/8 & 211 &  0.0 & 211.0 & discard \\
\bottomrule
\end{tabular}
\end{center}

Full MDL: 3 improvements kept, baseline $279.0 \to 97.0$,
\textbf{65.2\% reduction}. Running the same loop with $L(D \mid H)$ only:
0 improvements kept, 0\% reduction, stuck at baseline forever. The baseline
already passes all 8 tests---$L(D \mid H) = 0$. Without $L(H)$, the score
is at its minimum from step 0. The loop is completely blind.

The decisive case is \texttt{06\_overengineered}. It passes all 8 tests.
$L(D \mid H) = 0$. $L(H) = 211$. The full MDL loop discards it
immediately---it costs 114 bits more than the current best with no benefit.
The $L(D \mid H)$-only loop cannot distinguish it from
\texttt{05\_use\_builtin}. Both score 0.0.

This is the Karpathy gap. Karpathy's system minimizes \texttt{val\_bpb}:
validation bits per byte, literally $L(D \mid H)$. The missing term is
$L(H)$---the description length of the training script, the model
architecture, the configuration. Adding it completes the MDL objective.

The structure is identical to the replication crisis
(Section~\ref{sec:connections}): an apparent result (unchanged validation
loss, a significant p-value) conceals a real cost that is not being
accounted for. The fix is the same in both cases---add the missing term.
At the scale of real AutoResearch loops running hundreds of iterations, a
system without $L(H)$ will accumulate complexity whenever doing so does not
hurt validation loss. The hypothesis grows. No validation signal warns that
something is wrong. The full MDL objective is $L(H) + L(D \mid H)$.
Minimizing it honestly is the difference between approximating Solomonoff
induction and falling systematically short of it.